% Independent mathematical derivation and root continuation, 2026-09-23.
% CTO1--28, including 22a--22c; original source bytes retained privately.
% Original test topology, additive factors and full support labels retained.
\section{Global theta majorants and the quotient that retains both endpoints}
\label{sec:cc-theta-order-cofinality}

\paragraph{Original sources and scope of this derivation.}
The source is Alain Connes, Caterina Consani and Matilde Marcolli,
\href{https://arxiv.org/abs/math/0703392v1}{\emph{The Weil proof and the
geometry of the adeles class space}}, original author file
\nolinkurl{Yuri70v2.tex}, canonical source
\texttt{PUBUNIT-5DE2A2D44F8B1470AC701FDA}. Its SHA-256 is
\par\noindent
\nolinkurl{b1f694cb934dd4fef829287f1a55dd8de21c67a481e552453c2f7e8fe0ae6fb5}.
\par
For this calculation the original passages read were lines 494--589,
623--640, 688--718, 1371--1462, 1631--1751, 2181--2284,
and 2286--2435. The exact locators are \texttt{SCK},
\texttt{rangecVdef}, \texttt{varthetaH1V}, \texttt{varthetafconvol},
\texttt{rhoFourier}, \texttt{degcorrZf}, \texttt{codegcorrZf},\newline
\texttt{PositSect}, \texttt{ZfbulletZf}, \texttt{fubini1}--\texttt{fubini4},
\texttt{nothappen} and \texttt{degmodifyV}.
The adelic measure conventions are also those of Alain Connes,
\href{https://arxiv.org/abs/math/9811068v1}{\emph{Trace formula in
noncommutative geometry and the zeros of the Riemann zeta function}},
Section III, as retained in the accompanying ABR derivation.

The original polynomial Gaussian comes from CCM's
\texttt{degmodifyV}. The prior programme proof GTK1--9 in
\nolinkurl{GLOBAL_THETA_WEIL_KERNEL.tex} already proves its theta image
strictly positive, its strong decay, its full Mellin formula and its
two endpoint values. Those facts are inputs, not new results here.
The original test topology, convolution and involution proofs are SWQ1--15 in
\nolinkurl{CC_SUPPORTED_WEIL_QUOTIENT.tex}. These precise source
passages and programme statements were read for this derivation.
The new construction proves cofinal domination for the entire original
real test space, an actual adelic antecedent for each majorant, and the
exact resulting ordered quotient with degree and codegree retained.
The last section distinguishes this test-function order from the
geometric effectivity used by Weil and discussed by CCM.

\subsection{The original strong test space and the existing theta test}

For the rational field retain the actual product and measures
\begin{equation}\tag{CTO1}
\begin{gathered}
C_{\mathbb Q}=\mathbb R_{>0}\times K,\qquad
K=\widehat{\mathbb Z}^{\times},\qquad d^\times c=dy/y\,du,
\quad \int_Kdu=1,\\
B={\bf S}(C_{\mathbb Q})
 =\bigcap_{\beta\in\mathbb R}|\cdot|^\beta\mathcal S(C_{\mathbb Q}),
\qquad B_{\mathbb R}=\{g\in B:g\text{ is real-valued}\},\\
T=\mathcal S(\mathbb R)\otimes
 \bigotimes_p'(\mathcal S(\mathbb Q_p),\mathbf1_{\mathbb Z_p}),\qquad
T_0=\{\xi\in T:\xi(0)=0,\ \int_{\mathbb A}\xi=0\},\\
E\xi(y,u)=\sum_{q\in\mathbb Q^\times}\xi(qy,qu),\qquad
V=E(T_0),\quad V_{\mathbb R}=V\cap B_{\mathbb R}.
\end{gathered}
\end{equation}
CCM's definition \texttt{SCK} uses the intersection of all real
weights. The symbol $B$ here does not mean the larger space with
only ordinary Schwartz decay in $\log y$.
At a fixed compact open level in $K$ its original seminorms are
\begin{equation}\tag{CTO2}
p_{A,m,j}(g)=\sup_{t\in\mathbb R,\ u\in K}
 e^{A|t|}(1+|t|)^m
          |\partial_t^j g(e^t,u)|,
\qquad A,m,j\in\mathbb Z_{\ge0}.
\end{equation}
The full Bruhat topology is the inductive limit over these finite
compact-unit levels. Every test considered below remains at a
specified finite level; the constructed majorants are at the trivial
level. All additive integrals use Lebesgue measure at infinity and
$\operatorname{vol}(\mathbb Z_p)=1$ at each finite place.

Retain the original seed and all its factors:
\begin{equation}\tag{CTO3}
\begin{gathered}
\eta(x)=\pi x^2(\pi x^2-\tfrac32)e^{-\pi x^2},\qquad
\xi(x_\infty,x_f)=\eta(x_\infty)
                       \prod_p\mathbf1_{\mathbb Z_p}(x_p),\\
h(y)=E\xi(y,u)=2\sum_{n\ge1}\eta(ny),\qquad H(t)=h(e^t),\\
\xi(0)=0,\quad \int_{\mathbb A}\xi=0,
\quad h\in V_{\mathbb R},\quad h(y)>0,
\quad h(y)=y^{-1}h(y^{-1}).
\end{gathered}
\end{equation}
These are the proved GTK1--5 inputs. The additive zero integral is
the exact cancellation $3/4-3/4$ of the quartic and quadratic
Gaussian moments. The factor two in $E\xi$ retains both rational
signs, and the finite factors permit exactly the nonzero integers.
All derivatives of $H$ satisfy the strong seminorms CTO2.

An explicit lower bound on the interval used below is
\begin{equation}\tag{CTO4}
H(t)\ge\kappa:=2\pi(\pi-\tfrac32)
                  \exp\{-\pi e^{3/2}\}>0
\qquad(|t|\le3/4).
\end{equation}
For $0\le t\le3/4$, retain the $n=1$ term in CTO3; all other
terms are positive. In that term $e^{2t}\ge1$,
$\pi e^{2t}-3/2\ge\pi-3/2$, and
$e^{-\pi e^{2t}}\ge e^{-\pi e^{3/2}}$. Their product proves
the bound. For negative $t$, CTO3 gives $H(t)=e^{-t}H(-t)$,
and its extra positive factor is at least one.

\subsection{A global majorant in the original image}

Fix any $g\in B_{\mathbb R}$, including its full dependence on $K$.
Define the single compact bump and its original integral by
\begin{equation}\tag{CTO5}
\begin{gathered}
\psi(t)=
\begin{cases}
\exp\{-1/(1-16t^2)\},&|t|<1/4,\\
0,&|t|\ge1/4,
\end{cases}\qquad b_\psi=\int_{\mathbb R}\psi(t)\,dt>0,\\
M_n=\sup_{u\in K,\ |t-n|\le1/2}|g(e^t,u)|,\qquad
a_n=\frac{M_n+e^{-n^2}}{\kappa b_\psi}\quad(n\in\mathbb Z),\\
W(t)=\sum_{n\in\mathbb Z}a_n\psi(t-n),\qquad
w(y,u)=W(\log y).
\end{gathered}
\end{equation}
The compact supremum exists and is finite: at the original finite
unit level it is the maximum of finitely many continuous real
functions on a compact interval. Each $a_n$ is strictly positive.
For every $A\ge0$, CTO2 implies
\[
M_n\le p_{A,0,0}(g)
      \exp\{-A\max(|n|-1/2,0)\}.
\]
It follows, using an arbitrarily larger exponential weight, that
\begin{equation}\tag{CTO6}
\sum_{n\in\mathbb Z}a_n e^{A|n|}(1+|n|)^m<\infty
\quad\hbox{for every }A,m\ge0.
\end{equation}
For the added term $e^{-n^2}$ this follows directly from its
quadratic exponent. For the $M_n$ term choose the estimate with
weight $A+2$ and bound $(1+|n|)^m e^{-2|n|}$ by a summable
sequence. Noninteger $A,m$ are covered by larger integer values.

The bump sum is locally finite, and its supports are disjoint.
For every seminorm in CTO2, each derivative of its tail is bounded by
\begin{equation}\tag{CTO7}
p_{A,m,j}\left(\sum_{|n|>N}a_n\psi(\log y-n)\right)
\le\|\psi^{(j)}\|_\infty
\sup_{|n|>N}a_n e^{A(|n|+1/4)}(1+|n|+1/4)^m
\longrightarrow0.
\end{equation}
Thus $w\in B_{\mathbb R}$ and the displayed series converges in
the original topology at the trivial compact-unit level, with every
derivative. In particular all integrals
$\int|W(t)|e^{A|t|}(1+|t|)^m\,dt$ are finite, by integrating
the bumps and using CTO6. The coefficients do not need to dominate
the derivatives of $g$; they construct $w$ and its derivatives
from a fixed bump, while the weighted values of $g$ suffice for
the domination asserted below.

Use the original multiplicative convolution and keep its Haar factors:
\begin{equation}\tag{CTO8}
v(y,u)=(h*w)(y,u)
 =\int_K\int_0^\infty h(y/a)w(a,k)\,\frac{da}{a}\,dk
 =\int_{\mathbb R}H(\log y-t)W(t)\,dt.
\end{equation}
The compact integral has value factor one; both inputs are radial,
so the resulting function is independent of $u$.
Convolution remains in $B$ with its original seminorms. For
completeness the inequality
$|s|\le|t|+|s-t|$ and
$1+|s|\le(1+|t|)(1+|s-t|)$ gives
\[
p_{A,m,j}(h*w)
 \le2p_{A+1,m,0}(w)p_{A,m,j}(h),
\]
because $\int_{\mathbb R}e^{-|t|}dt=2$ and $\int_Kdk=1$.
The same integrable majorant justifies every derivative and proves
the strong convergence after replacing $W$ by its finite sums.

Given a real $s$, choose an integer $n$ with $|s-n|\le1/2$.
For every $t$ in the support of $\psi(t-n)$, one has
$|s-t|\le3/4$. Since all terms in CTO8 are nonnegative,
CTO4 gives the exact domination
\begin{equation}\tag{CTO9}
v(e^s,u)\ge\kappa a_n b_\psi
 =M_n+e^{-n^2}>|g(e^s,u)|
\qquad(s\in\mathbb R,\ u\in K).
\end{equation}
This is a global inequality on all of $C_{\mathbb Q}$, not a
finite interval or a finite set of compact-unit characters.

\subsection{The actual adelic antecedent and all convergence}

Membership in the algebraic image $V$, rather than merely its closure,
has an explicit original Schwartz antecedent:
\begin{equation}\tag{CTO10}
\eta_W(x)=\int_{\mathbb R}W(t)\eta(e^{-t}x)\,dt,
\qquad
\xi_W(x_\infty,x_f)=\eta_W(x_\infty)
                          \prod_p\mathbf1_{\mathbb Z_p}(x_p).
\end{equation}
For $P_{M,j}(\eta)=\sup_x(1+|x|)^M|\eta^{(j)}(x)|$, a real
dilation has the precise seminorm bound
\begin{equation}\tag{CTO11}
\begin{split}
\sup_x(1+|x|)^M
 \left|\frac{d^j}{dx^j}\eta(e^{-t}x)\right|
&\le P_{M,j}(\eta)e^{M\max(t,0)-jt}\\
&\le P_{M,j}(\eta)e^{(M+j)|t|}.
\end{split}
\end{equation}
The first inequality uses the actual factor $e^{-jt}$ and
$(1+e^t|z|)^M\le e^{M\max(t,0)}(1+|z|)^M$.
CTO6--7 make the integral of this bound against $|W(t)|$ finite
in every seminorm. Differentiating on compact $x$ intervals and
then applying the same bound proves that CTO10 is Schwartz and
that its integral converges in the Schwartz topology. One may
equivalently first integrate each compact bump and then sum;
the resulting series is absolutely convergent in each seminorm.
This constructs one real Schwartz function and a single permissible
adelic tensor with every finite reference factor retained. It
does not introduce an infinite algebraic tensor expansion.

Its two source endpoints vanish with absolute justifications:
\begin{equation}\tag{CTO12}
\begin{gathered}
\xi_W(0)=\int W(t)\eta(0)\,dt=0,\\
\int_{\mathbb A}\xi_W
 =\left(\prod_p\operatorname{vol}(\mathbb Z_p)\right)
    \int_{\mathbb R}W(t)e^t\,dt\int_{\mathbb R}\eta(x)\,dx
 =1\cdot\int_{\mathbb R}W(t)e^t\,dt\cdot(3/4-3/4)=0.
\end{gathered}
\end{equation}
For the second exchange the integral of the absolute values is
$\int|W(t)|e^tdt\int|\eta(x)|dx<\infty$. Thus
$\xi_W\in T_0$ in the original topology.

To prove its exact theta image, put
$P_2=\sup_x(1+|x|)^2|\eta(x)|$. For every $z>0$, monotone
integral comparison gives
\[
\sum_{n\ge1}|\eta(nz)|
 \le P_2\sum_{n\ge1}(1+nz)^{-2}
 \le P_2\int_0^\infty(1+uz)^{-2}du=P_2/z.
\]
Consequently, at each $y>0$,
\begin{equation}\tag{CTO13}
\begin{gathered}
\sum_{n\ne0}\int_{\mathbb R}
 |W(t)\eta(ne^{-t}y)|\,dt
 \le\frac{2P_2}{y}\int|W(t)|e^t\,dt<\infty,\\
E\xi_W(y,u)
 =\int_{\mathbb R}W(t)
                  2\sum_{n\ge1}\eta(ne^{-t}y)\,dt
 =v(y,u),\qquad v\in V_{\mathbb R}.
\end{gathered}
\end{equation}
The finite argument is still $nu\in\widehat{\mathbb Z}$ precisely
for integral $n$; unit coordinates and the two signs have not
been dropped. CTO13 uses an absolutely convergent exchange for
this fixed $y$ and $t$ integral. It does not interchange a
multiplicative endpoint integral with the original integer sum,
the invalid exchange in CCM's \texttt{nothappen}.

\subsection{The original Mellin factors and the endpoint cost}

Use CCM's Fourier convention without an inverse character:
\begin{equation}\tag{CTO14}
\widehat f(\chi,s)
 =\int_K\int_0^\infty f(y,u)\chi(u)y^s\,\frac{dy}{y}\,du,
\qquad M_f(s)=\widehat f(1,s).
\end{equation}
Every such transform is entire, with derivatives obtained by
inserting $(\log y)^j$, by CTO2. Set
\[
B_\psi(s)=\int_{-1/4}^{1/4}\psi(t)e^{st}\,dt,
\qquad
M_w(s)=\int_{\mathbb R}W(t)e^{st}\,dt
       =B_\psi(s)\sum_{n\in\mathbb Z}a_ne^{sn}.
\]
CTO6 proves local uniform convergence of this series and of all
its derivatives. Absolute weighted Fubini in CTO8 gives the full
original arithmetic formula
\begin{equation}\tag{CTO15}
\begin{split}
M_v(s)&=M_h(s)M_w(s),\\
M_h(s)&=2\zeta(s)\left\{
 \frac{\pi^2\Gamma((s+4)/2)}{2\pi^{(s+4)/2}}
 -\frac{3\pi\Gamma((s+2)/2)}{4\pi^{(s+2)/2}}
 \right\}\qquad(\operatorname{Re}s>1).
\end{split}
\end{equation}
The first identity holds for every complex $s$. The second is
GTK7, obtained on its specified domain by integrating the two
original Gaussian monomials separately. Its equality to the
entire Mellin integral specifies continuation at each removable
point of the product; no individual zero, pole, factor or summand
is deleted. In particular the original zeta function has not been
replaced by a completed one.

The known original endpoint values are $M_h(0)=M_h(1)=1/4$
(GTK9). Therefore the exact costs of the majorant are
\begin{equation}\tag{CTO16}
\begin{gathered}
C_v:=M_v(0)=\tfrac14\int_{\mathbb R}W(t)\,dt
          =\tfrac14 b_\psi\sum_na_n>0,\\
D_v:=M_v(1)=\tfrac14\int_{\mathbb R}W(t)e^t\,dt
          =\tfrac14 B_\psi(1)\sum_na_ne^n>0.
\end{gathered}
\end{equation}
These are codegree and degree respectively in CCM's
\texttt{codegcorrZf} and \texttt{degcorrZf}. They coexist with the
two additive zeros CTO12 because their receiving integrals are
different, and CTO13 does not perform the invalid endpoint Fubini
exchange. Adding this particular $v$ changes both endpoint values
by the displayed positive amounts.

The construction retains the full spectral jets. For a nontrivial
zero $\rho$ of the original $\zeta$ of multiplicity $m_\rho$, the
Gamma expressions in CTO15 are holomorphic at $\rho$ and $M_w$
is entire. Thus
\begin{equation}\tag{CTO17}
\begin{gathered}
M_v^{(j)}(\rho)=0\quad(0\le j<m_\rho),\\
\widehat v(\chi,s)=0\quad\text{for every nontrivial continuous
 character }\chi\text{ of }K\text{ and every }s.
\end{gathered}
\end{equation}
The first statement retains at least the full original multiplicity;
additional zeros of $M_w$ may increase it. The second follows from
radiality and Haar invariance: if $\chi(u_0)\ne1$, then
$\int_K\chi(u)du=\chi(u_0)\int_K\chi(u)du$, forcing that integral
to zero. All its $s$ derivatives are zero as well. Thus adding the
constructed majorant leaves those values and jets unchanged.

\subsection{The complete quotient-cone consequence and fixed endpoints}

Let $B_{\mathbb R,+}$ be the cone of everywhere nonnegative
real tests, and let $q:B_{\mathbb R}\to
Q_{\mathbb R}:=B_{\mathbb R}/V_{\mathbb R}$ be the algebraic
quotient. CTO9 and CTO13 prove, for every real $g$, the two
representatives
\begin{equation}\tag{CTO18}
g+v>0,\qquad g-v<0,\qquad q(g+v)=q(g-v)=q(g).
\end{equation}
In particular
\[
q(B_{\mathbb R,+})=Q_{\mathbb R},\qquad
q(\{f\in B_{\mathbb R}:f>0\})=Q_{\mathbb R}.
\]
This evaluates the entire image cone. Its associated preorder is
the universal relation: every quotient class is greater than or
equal to every other one. It does not separate any distinct
quotient classes. The same conclusion holds after replacing
$V$ by its closure, because the very same actual antecedents
already lie in $V$; no new closed-image assertion is needed.

Nonzero nonnegative tests have strictly positive values at both
endpoints. Indeed continuity at a positive value gives an open
set of positive Haar measure on which the function is positive;
both $du\,dy/y$ and $y\,du\,dy/y$ assign positive measure there.
Their integrals are finite by CTO2. Therefore, putting
$B_{\mathbb R,00}=\ker(M_0,M_1)$, one has
\begin{equation}\tag{CTO19}
B_{\mathbb R,+}\cap\ker M_0
=B_{\mathbb R,+}\cap\ker M_1
=B_{\mathbb R,+}\cap B_{\mathbb R,00}=\{0\}.
\end{equation}
Consequently a nonzero nonnegative correction cannot preserve
either zero endpoint change, and in particular cannot preserve
both prescribed degree and codegree.

There is a precise necessary bound for any nonnegative
representative with fixed endpoints $C=M_f(0)$, $D=M_f(1)$.
For every continuous character $\chi$ of $K$ and every
$\rho=\sigma+i\gamma$ with $0<\sigma<1$, the original integral
and H\"older's inequality give
\begin{equation}\tag{CTO20}
\begin{split}
|\widehat f(\chi,\rho)|
&\le\int_K\int_0^\infty f(y,u)y^\sigma\,\frac{dy}{y}\,du\\
&=\int_K\int_0^\infty
       f(y,u)^{1-\sigma}\{yf(y,u)\}^{\sigma}
                                          \,\frac{dy}{y}\,du\\
&\le C^{1-\sigma}D^\sigma.
\end{split}
\end{equation}
Here $|\chi(u)|=|y^{i\gamma}|=1$, and the H\"older exponents
are $1/(1-\sigma)$ and $1/\sigma$ on the original measure.
For the zero test both sides are zero; otherwise CTO19 proves
$C,D>0$. This is a receiving inequality with an endpoint budget.
CTO9 by itself supplies no uniform values of that budget, while
CTO16 gives its exact cost for each input and CTO17 retains
the indicated spectral data.

\subsection{An actual ordered extension retaining both endpoints}

The failed quotient order defines a more informative exact receiver.
Set
\begin{equation}\tag{CTO21}
N=V_{\mathbb R}\cap B_{\mathbb R,00},\qquad
\mathcal E_{\rm ord}=B_{\mathbb R}/N,\qquad
\mathcal E_{{\rm ord},+}=\pi_N(B_{\mathbb R,+}).
\end{equation}
These are algebraic quotients with the original topology available;
no closedness claim about $N$ or the cone is made. The real linear
endpoint map $\epsilon(f)=(M_f(0),M_f(1))$ descends to
$\mathcal E_{\rm ord}$.
For $b>0$ write $h_b(y)=h(y/b)$, with the exact source
$\eta(x_\infty/b)\prod_p\mathbf1_{\mathbb Z_p}$.
The source still has evaluation zero and additive integral
$b(3/4-3/4)=0$, so $h_b\in V_{\mathbb R}$.
The multiplicative substitutions give
$\epsilon(h_b)=(1/4,b/4)$. In particular the explicit linear map
\[
j(C,D)=\bigl[4(2C-D)h+4(D-C)h_2\bigr]_N
\]
has endpoints $(C,D)$. A different representative in $V_{\mathbb R}$
with the same endpoints differs by exactly an element of $N$,
so this is the canonical identification of the kernel with its
endpoint pair, independent of that selected formula. We obtain
the exact sequence
\begin{equation}\tag{CTO22}
0\longrightarrow\mathbb R^2
 \xrightarrow{\ j\ }\mathcal E_{\rm ord}
 \xrightarrow{\ q_{\rm ord}\ }Q_{\mathbb R}
 \longrightarrow0,
\qquad \epsilon j=I_{\mathbb R^2}.
\end{equation}
Indeed the middle kernel is $V_{\mathbb R}/N$, whose endpoint
map is injective by the definition of $N$ and onto by the
displayed $h,h_2$ formula. Surjectivity of $q_{\rm ord}$ is
the original quotient map. Every map in this sequence has
therefore been specified on its original representatives.

This linear extension splits canonically; its additional information
is its positive cone, not an asserted non-split linear structure.
The exact isomorphism and inverse are
\begin{equation}\tag{CTO22a}
\begin{gathered}
\Upsilon:\mathcal E_{\rm ord}\xrightarrow{\ \cong\ }
          Q_{\mathbb R}\oplus\mathbb R^2,\qquad
\Upsilon([f]_N)=([f]_V,M_f(0),M_f(1)),\\
\Upsilon^{-1}(q,C,D)
=[g]_N+j(C-M_g(0),D-M_g(1)),\qquad [g]_V=q.
\end{gathered}
\end{equation}
If $g$ is replaced by $g+v$ with $v\in V_{\mathbb R}$, then
$[v]_N=j(M_v(0),M_v(1))$ by the exact kernel identification.
The two added terms therefore cancel in the inverse formula.
Thus the inverse is independent of the representative and is
linear. Applying either composite proves the identity, so
surjectivity and injectivity are both established. The real
linear section of $q_{\rm ord}$ is $q\mapsto\Upsilon^{-1}(q,0,0)$.
No degree or codegree contribution is omitted in this decomposition.

The original convolution algebra structure is retained too. To recall
its required source calculation, for $f\in B$ and $\xi\in T_0$ the
actual section $C_{\mathbb Q}\to\mathbb A^\times$ of CTO1 gives
\[
\xi_f(x)=\int_C f(a)\xi(a^{-1}x)\,d^\times a,\qquad
E\xi_f=f*E\xi.
\]
SWQ11--14 proves this directly. The compact unit orbit of a fixed
finite Bruhat test spans a finite-dimensional space, and each real
Schwartz seminorm is bounded by a constant times
$\int_C|f(a)|\max(1,|a|)^M|a|^{-j}d^\times a<\infty$.
Its endpoints are
$\xi_f(0)=\xi(0)M_f(0)$ and
$\int\xi_f=(\int\xi)M_f(1)$, both zero. For each fixed receiving
idele the absolute rational sum has a majorant
$C_\xi(1+|a|)$; integrating against $|f(a)|$ justifies the
last displayed identity. These are the original source maps
and factors, proving $V$ is a convolution ideal. On real tests,
complex conjugation shows the same for $V_{\mathbb R}$.
Absolute multiplicative Fubini gives
$M_{f*g}(0)=M_f(0)M_g(0)$ and
$M_{f*g}(1)=M_f(1)M_g(1)$, retaining $|ab|=|a||b|$ in the
second identity. Thus $N$ is a convolution ideal and
\begin{equation}\tag{CTO22b}
\begin{gathered}
\Upsilon([f*g]_N)
 =([f]_V*[g]_V,M_f(0)M_g(0),M_f(1)M_g(1)),\\
j(C,D)*[f]_N=j(CM_f(0),DM_f(1)).
\end{gathered}
\end{equation}
This proves the direct-product algebra structure on the same
canonical coordinates; it introduces no convolution identity
among the original test functions.

The original involution is
$f^\sharp(c)=|c|^{-1}\overline{f(c^{-1})}$.
With the self-dual additive Fourier transform retained by SWQ6--10,
the full Poisson endpoint comparison is
\begin{equation}\tag{CTO22c}
\begin{gathered}
(E\xi)^\sharp(c)
 =E\widehat{\overline\xi}(c)
     -|c|^{-1}\overline{\xi(0)}
     +\overline{\int_{\mathbb A}\xi},\\
M_{f^\sharp}(0)=\overline{M_f(1)},\qquad
M_{f^\sharp}(1)=\overline{M_f(0)}.
\end{gathered}
\end{equation}
The first equality retains the original real Fourier kernel
$e^{-2\pi ixv}$, finite conductor $\mathbb Z_p$ and both
source endpoints; SWQ7--10 proves it by Poisson summation.
The Fourier antecedent has endpoints
$\overline{\int\xi},\overline{\xi(0)}$, so $V^\sharp=V$.
The last two equations follow by $c\mapsto c^{-1}$ in the
original Haar integrals, including the factor $|c|^{-1}$.
Thus sharp preserves $N$ and, on the real ordered extension,
acts by the original quotient involution and interchange of
the two real endpoints. It preserves its positive cone because
the multiplying modulus is positive. No Hilbert-space positivity
is inferred from this algebraic involution.

The cone of CTO21 is pointed and generating. If $x$ and $-x$
have nonnegative representatives $f$ and $g$, then $f+g\in N$,
so $M_{f+g}(0)=0$. CTO19 implies $f+g=0$, hence $f=g=0$
and $x=0$. Generating follows from CTO18:
$[g]_N=[g+v]_N-[v]_N$, with both terms nonnegative.
Its projection to $Q_{\mathbb R}$ is the entire quotient by
CTO18. Its intersection with the exact kernel in CTO22 is
completely determined:
\begin{equation}\tag{CTO23}
j^{-1}\bigl(\mathcal E_{{\rm ord},+}\cap\ker q_{\rm ord}\bigr)
=\{(C,D):C>0,\ D>0\}\cup\{(0,0)\}.
\end{equation}
For necessity a representative of a kernel class belongs to
$V_{\mathbb R}$, and CTO19 gives both strict inequalities unless
it is zero. For sufficiency, if $C,D>0$, the actual function
$4C h_{D/C}>0$ belongs to $V_{\mathbb R}$ and has endpoints
$4C(1/4,(D/C)/4)=(C,D)$. This gives its exact representative
in the image of $j$. The origin is represented by zero.
Thus the two coordinate axes, apart from their origin, are not
silently included in this cone.
Under CTO22a this cone is therefore not the naive product
$Q_{\mathbb R}\times\mathbb R_{\ge0}^2$: for example its
kernel coordinate $(0,1,0)$ has no nonnegative representative,
by CTO19 and CTO23. The calculation of the kernel cone does
not assert that the complete positive cone at each prescribed
nonzero spectral class and each prescribed endpoint pair has
been evaluated; CTO20 supplies an exact constraint there.

For an idele class $a$, the original translation
$U_af(c)=f(a^{-1}c)$ preserves pointwise nonnegativity, $V$,
and $N$. Invariance of $V$ follows already at the source by
$EU_a=U_aE$, and the additive source integral is multiplied
by $|a|$ while its zero remains zero. The original multiplicative
change of variable gives
$\epsilon(U_af)=(M_f(0),|a|M_f(1))$.
Consequently CTO22 is equivariant with kernel action
$(C,D)\mapsto(C,|a|D)$, and the cone CTO23 is invariant.
The real linear maps in the sequence, rather than a selected
nonlinear majorant, carry this equivariance.

\subsection{The closed arithmetic quotient and the exact topology of the cone}

The algebraic calculation has an exact map to CCM's quotient by the
closed theta image. Keep the original topology CTO2 and define
\begin{equation}\tag{CTO25}
\begin{gathered}
s(C,D)=4(2C-D)h+4(D-C)h_2,\qquad P=I-s\epsilon,\\
\overline Q_{\mathbb R}=B_{\mathbb R}/\overline{V_{\mathbb R}},\qquad
\overline{\mathcal E}_{\rm ord}=B_{\mathbb R}/\overline N.
\end{gathered}
\end{equation}
The bar on $\overline{\mathcal E}_{\rm ord}$ here denotes this
specified quotient, not an unproved completion of an ordered space.
The endpoint integrals are continuous: each is bounded by a constant
times $p_{2,0,0}$ on each compact-unit level. The fixed two-function
map $s$ is continuous, has image in $V_{\mathbb R}$, and satisfies
$\epsilon s=I$. Hence $P$ is a continuous projection onto
$\ker\epsilon$, with $P(V_{\mathbb R})=N$.
Continuity of $\epsilon$ first gives
$\overline N\subseteq\overline{V_{\mathbb R}}\cap\ker\epsilon$.
Conversely, if $z$ is in the latter intersection, a net $v_i\in
V_{\mathbb R}$ tends to $z$ and $Pv_i\in N$ tends to $Pz=z$.
Therefore
\begin{equation}\tag{CTO26}
\overline N=\overline{V_{\mathbb R}}\cap\ker\epsilon,\qquad
\overline\Upsilon:\overline{\mathcal E}_{\rm ord}
\xrightarrow{\ \cong\ }
\overline Q_{\mathbb R}\oplus\mathbb R^2,\qquad
[f]_{\overline N}\longmapsto
([f]_{\overline{V_{\mathbb R}}},M_f(0),M_f(1)).
\end{equation}
This is a topological real vector-space isomorphism, not only an
algebraic bijection. Indeed $f\mapsto(Pf,\epsilon f)$ is a
continuous isomorphism from $B_{\mathbb R}$ to
$\ker\epsilon\oplus\mathbb R^2$, with continuous inverse
$(z,C,D)\mapsto z+s(C,D)$. It takes $V_{\mathbb R}$ to
$N\oplus\mathbb R^2$ and its closure to
$\overline N\oplus\mathbb R^2$. Passing to these quotient
topologies proves (CTO26) and its inverse. The same argument
without bars also makes CTO22a a topological isomorphism in the
possibly non-Hausdorff quotient topologies.

Complex conjugation preserves $V$, since it commutes with $E$.
Taking real parts of a convergent net proves
$\overline{V_{\mathbb R}}=\overline V\cap B_{\mathbb R}$.
Thus (CTO26) uses the actual real part of the closed-image quotient.
Multiplication by a fixed original test is continuous, by the
weighted convolution estimate following CTO8. Consequently the
closures of the ideals $V_{\mathbb R}$ and $N$ remain ideals.
The endpoint characters give exactly the product law CTO22b on
(CTO26); it is an isomorphism of the original convolution algebras.
Sharp and every fixed idele translation are continuous by their
explicit formulas, so CTO22c and its original actions also descend.
No closedness assumption on $V$ or $N$ has been inserted.

Give $\overline{\mathcal E}_{\rm ord}$ the image of
$B_{\mathbb R,+}$, without closing that cone. Its cone is pointed
by the identical CTO19 argument, since every element of
$\overline N$ has both endpoints zero. It is generating, its
projection to $\overline Q_{\mathbb R}$ is onto, and its exact
kernel cone is still CTO23: necessity uses the positive endpoint
integrals and sufficiency uses the same function $4Ch_{D/C}$.

Both the algebraic and the closed-image versions have a positive
cone which is \emph{not closed} in their stated quotient topology.
For the first, and identically for the second, continuity of $j$ gives
\begin{equation}\tag{CTO27}
j(1,1/n)=[4h_{1/n}]_N\in\mathcal E_{{\rm ord},+},\qquad
j(1,1/n)\longrightarrow j(1,0)
\notin\mathcal E_{{\rm ord},+}.
\end{equation}
The convergence is of quotient classes through the fixed continuous
map $s$; convergence of $4h_{1/n}$ in the original test space is
not asserted. The excluded limit follows from CTO23. Similarly
$j(0,1)\not\ge0$, whereas
$j(0,1)+\varepsilon j(1,1)=j(\varepsilon,1+\varepsilon)>0$ for
every $\varepsilon>0$. Thus this order is not Archimedean: positive
perturbations by arbitrarily small multiples of this order unit do
not force positivity at their limit.

There is nevertheless an actual order unit $u_{\rm ord}=j(1,1)$
for the entire extension. For $x=[g]_N$, the constructed $v>|g|$
has endpoints $C_v,D_v>0$ and $[v]_N=j(C_v,D_v)$. For every
real $R>\max(C_v,D_v)$,
\begin{equation}\tag{CTO28}
Ru_{\rm ord}\pm x
=[v\pm g]_N+j(R-C_v,R-D_v)
\in\mathcal E_{{\rm ord},+}.
\end{equation}
Both summands have actual positive representatives, by CTO9 and
CTO23. This proves the order-unit assertion on every class and
also on (CTO26), keeping its full individual cost. It does not
justify replacing this cone by its closure, passing a positive
limit back into it, or declaring the original arithmetic trace
positive on it. Those inferences would contradict or exceed the
computed cone properties.

\subsection{Geometric scope and every original support label}

CCM's source discussion \texttt{PositSect} uses Riemann--Roch
on a curve: after adjusting a correspondence's degree, a principal
correspondence produces geometric effectivity in its linear
equivalence class. The preceding definition of effective
correspondence is a sum of curves with nonnegative integer
multiplicities, with intersection products and Frobenius action.
The later \texttt{degmodifyV} gives freedom to adjust degree by
an actual trivial theta correspondence. The calculation above
retains and strengthens the test-function part of that comparison:
it gives a positive representative in every real test quotient
class, computes its endpoint costs, and constructs the exact
ordered extension retaining them. It does not identify a
pointwise-nonnegative test function with an effective divisor,
construct the missing principal-divisor relation, or establish
the intersection inequality \texttt{ZfbulletZf}. The universal
cone in CTO18 is consequently an evaluated defect of this
particular proposed order, not a no-go statement about all
geometric notions of effectivity.

For an arbitrary bounded distributive lattice $L$ retain
\begin{equation}\tag{CTO24}
G_L(M)=\{(0,\lambda):\lambda\in L\}\cup(M\times\{1_L\}),
\qquad e=(0,1_L),\quad\tau=(0,0_L),\qquad
A^\uparrow(z,\lambda)=(Az,\lambda)
\end{equation}
for every displayed linear map $A$ on its specified amplitude
spaces. The originating support source, by The Clankers, is
\href{https://github.com/KokunoYumeto/zeta-function-research-reader/blob/a2851f9eb352e7e4376bc629de86fa4ed9c1c434/workbenches/splitzero-tandem/foundations/split-support-geometry-v11/split_support_geometry_arithmetic_curve_v11.tex}{\emph{Split Support Geometry}, v11,
Definition \texttt{def:lattice-split}}.
Addition retains join and scalar multiplication retains meet.
All quotients above use the same-label amplitude congruence;
CTO22 lifts as the corresponding exact amplitude description
under those congruences. Every zero amplitude retains its
assigned label, so the positive theta majorant at top support
maps to $e$ in $G_L(Q_{\mathbb R})$, not to $\tau$.

The cofinal choice $g\mapsto v$ is nonlinear and is asserted
as a strict majorant construction for top-supported inputs.
A lower-supported zero $(0,\lambda)$ has no admissible nonzero
amplitude at that label. For it choose $(0,\lambda)$ itself as
a non-strict majorant. Thus no strict positive amplitude is
inserted into a forbidden lower support stratum. The convergent
series and integrals CTO5--13 are operations on amplitude
functions with one admissible unchanged label; no infinite
join or integration on an arbitrary lattice is postulated.
The endpoint pair in CTO22 likewise has one common label,
not independently assigned labels on its two coordinates.

This constructs a pointed ordered extension with exact degree
and codegree, while evaluating why the pointwise cone on the
unextended spectral quotient fails to separate its classes.
It supplies neither an assumed positivity theorem for the
original Weil form nor an RH conclusion.
